\begin{helper}[Missing colors under a two-color component interchange]
Let \(\varphi\) be a proper partial edge-coloring, let
\(\alpha\ne\beta\), let \(C\) be one
\((\alpha,\beta)\)-component of the two-color incidence graph, and let
\(\varphi'\) be obtained by interchanging \(\alpha\) and \(\beta\) on
\(C\).  Then the following hold.
\begin{enumerate}
\item If a vertex \(v\) is outside \(C\), then its entire missing-color
set is unchanged.
\item If \(v\in C\) misses \(\alpha\) under \(\varphi\), then it misses
\(\beta\) under \(\varphi'\); symmetrically, if it misses \(\beta\)
under \(\varphi\), then it misses \(\alpha\) under \(\varphi'\).
\item At every vertex, missingness of every color outside
\(\{\alpha,\beta\}\) is unchanged.
\end{enumerate}
\end{helper}
\begin{proof}
Write \(H\) for the incidence graph
\noderef{TwoColorIncidenceGraph} and use the component
\noderef{TwoColorComponent}.  Fix a vertex \(v\) and an incident edge
\(e\).  If the edge-node \(e\) belongs to \(C\) and its old value is
\(\alpha\) or \(\beta\), then the adjacency rule in
\noderef{TwoColorIncidenceGraph} joins \(e\) to the vertex-node \(v\),
so \(v\in C\).  If instead the edge-node belongs to \(C\) but its old
value is neither affected color (including the possibility that it is
uncolored), both nested color tests in
\noderef{TwoColorComponentSwap} fail and the value of \(e\) is fixed.
This second alternative includes the isolated edge-root case, in which
an off-color or uncolored edge-node can by itself be the selected
component.  Finally, if the edge-node is outside \(C\), the outer
component test in \noderef{TwoColorComponentSwap} fixes its value.

Conversely, if \(v\in C\) and an incident edge has old color
\(\alpha\) or \(\beta\), the same adjacency rule joins its edge-node to
\(v\), and hence that edge-node lies in \(C\).

Suppose first that \(v\notin C\).  The first observation shows that no
incident edge whose edge-node lies in \(C\) can have old value
\(\alpha\) or \(\beta\), since that would force \(v\in C\).  Every such
edge is therefore fixed by the nested tests, and every incident edge
outside \(C\) is fixed by the outer test.  Thus every exact value on an
edge incident with \(v\) is fixed.  By the definitions of present and
missing colors
\noderef{PresentColors} and \noderef{MissingColors}, the whole
missing-color set at \(v\) is unchanged.

Now suppose \(v\in C\) and \(v\) misses \(\alpha\) under \(\varphi\).
If an incident edge had color \(\beta\) after the interchange, then
either its edge-node was outside \(C\), in which case
\noderef{TwoColorComponentSwap} says that it already had color
\(\beta\), or its edge-node was inside \(C\).  The first case is
impossible: that old \(\beta\)-edge is adjacent to \(v\) in
\noderef{TwoColorIncidenceGraph}, so the edge-node belongs to the same
component as \(v\).  In the second case, the three nested branches in
\noderef{TwoColorComponentSwap}, together with
\(\alpha\ne\beta\), show that a new \(\beta\)-value can only have come
from an old \(\alpha\)-value.  That contradicts \(v\)'s old
missing-\(\alpha\) status by \noderef{MissingColors}.  Hence no
incident edge has color \(\beta\) afterward, so \(v\) misses
\(\beta\).  Interchanging the roles of \(\alpha\) and \(\beta\) proves
the symmetric assertion.  The assumed properness is the condition under
which these components are the usual two-color chains, as expressed by
\noderef{IsProperEdgeColoring}; the pointwise exclusions above prove the
claimed missing-color implications.

Finally, \noderef{TwoColorComponentSwap} changes only values equal to
\(\alpha\) or \(\beta\).  More explicitly, an old value
\(c\notin\{\alpha,\beta\}\) fails both nested color tests and is fixed;
conversely, neither affected-color branch can produce \(c\), so a new
\(c\)-value must be the fixed old \(c\)-value.  Thus an incident edge
has color \(c\) before the interchange if and only if it has color
\(c\) afterward.  Applying \noderef{PresentColors} and then taking
complements via \noderef{MissingColors} proves the last assertion.
\end{proof}
